\documentclass[11pt]{article}
\usepackage[margin=1in]{geometry}
\usepackage{amsmath,amssymb,amsthm,hyperref}
\newtheorem{theorem}{Theorem}
\newtheorem{proposition}[theorem]{Proposition}
\newtheorem{lemma}[theorem]{Lemma}
\newtheorem{corollary}[theorem]{Corollary}
\newcommand{\F}{\mathbb F}
\newcommand{\rank}{\operatorname{rank}}
\newcommand{\dist}{\operatorname{dist}}
\title{Lower bounds for time-dependent layout measurements}
\author{Chris Blaschke\\Ax0s}
\date{October 31, 2026; revised November 2, 2026}
\begin{document}
\maketitle
\begin{abstract}
We prove lower bounds on the accounting exponent of OpenAI's layout interface for bounded-incidence, balanced-cut expanding scopes. The proof couples arbitrary time-dependent measurement spaces through a determinant. Under a capacity condition the charged exponent is at least a constant times $(\log\log n)^{3/2}/\sqrt{\log n}$. We also prove a small-measurement bound in auxiliary coordinates and a separate partner-time inequality. The partner times cannot be dropped: Clifford-algebra polynomial matrices refute the far-pair matrix bound conjectured in our first version, making the measured rank an arbitrarily small fraction of $m(\Delta-R)$ even over $\mathbb Q$. These are not general algorithmic running-time lower bounds.
\end{abstract}
\section{Introduction}
OpenAI~\cite{openai} and Claude (Anthropic)~\cite{anthropic} independently refute the exponential time hypothesis, with counting algorithms of running time $2^{n/(\log n)^{1/2+o(1)}}$ and $2^{O(n\log\log n/\sqrt{\log n})}$ respectively. We ask whether the square-root logarithm is forced by the accounting for OpenAI's layout interface, rather than by the problem itself.

Tulsyan~\cite{tulsyan} proves degree barriers for the fixed-$M_0$ template, in which $M_0$ has the prescribed Hamming support and fixes the kernel of a common measurement matrix. On the full label space, his first bound yields $R\ge(1-o(1))\sqrt{Q\ln(1/\varepsilon)}$ in the even-alphabet moderate-deviation regime; his second is capacity-free and uses the Arora--Rao--Vazirani theorem~\cite{arv}. He asks how to couple corrections that vary with time. Olivier~\cite{olivier} independently proves a degree lower bound for each fixed matrix satisfying OpenAI's equation (4.4). Neither fixed-$M_0$ result settles general time-dependent layouts.

Our determinant inequality couples all evaluation points without identifying their correction spaces. On expanding scopes, it yields lower bounds on the layout accounting exponent, and on the full auxiliary label space it gives a small-measurement degree bound. A separate interpolation argument proves a partner-time inequality and an exact result for a family of binary coset instances. Neither argument establishes a capacity-free barrier for general layouts. We show that partner times are necessary for the second: when each direction has only one symmetric time, the far-pair matrix bound conjectured in our first version fails (Theorem~\ref{thm:far}).

\section{The layout interface}
Let $[n]$ be the variables and $I_f\subseteq[n]$ the scope of each factor $f$. OpenAI's layout has integers $r\ge2$, $D\ge1$, $1\le R<Q$ and $0<\tau\le1/2$. For each $t\in[Q]$ it specifies a partition $[n]=\bigsqcup_{g=1}^r S_{t,g}$, empty parts permitted, with $|S_{t,g}|\le2n/r$. Write $a_v(t)=g$ when $v\in S_{t,g}$, and write $d_f(t)$ for the assigned group of factor $f$.

For each $(t,g)$ there is an integer matrix $C_{t,g}$ on columns $S_{t,g}$, with $c_{t,g}$ rows, each of $\ell_1$ norm at most $D$, and $\sum_g c_{t,g}\le2\tau n$. For each incidence $v\in I_f$ there is a row $P_{f,v}(T)\in\mathbb Q[T]^{1\times n}$ of degree at most $R$ with $P_{f,v}(0)=e_v$. For every $g\ne d_f(t)$ and $\xi\in\ker C_{t,g}$, the row $P_{f,v}(t)$ annihilates the extension of $\xi$ by zero. Equivalently, its restriction to $S_{t,g}$ belongs to the row space of $C_{t,g}$. Finally, with $d_{\rm den}=(R-1)!Q!$, each coefficient of $d_{\rm den}P_{f,v}$ is an integer of $O(Q\log Q+\log D)$ bits, with an absolute constant. Our lower bounds do not use this coefficient condition.

For total product degree $d=O(n)$, the counting analysis of \cite{openai} gives the following upper-bound proxy for the normalized exponential cost, apart from polynomial factors:
\begin{equation}\label{eq:F}
 F=\frac RQ\log_2\frac{2Q}{R}+\tau\log_2(2D+1)+\frac1r.
\end{equation}
A lower bound on $F$ limits this particular accounting; it is not a lower bound on the running time of every algorithm, or even on an alternative implementation of the counting step.

\section{A determinant inequality}
The following result uses only the pointwise row-space condition, not expansion or integrality. It holds over any field $\F$ of characteristic zero, with distinct nonzero evaluation points $t$; signs are ordinary real signs applied to schedules.
\begin{theorem}\label{thm:det}
Choose signs $\epsilon_{t,g}\in\{-1,1\}$ and put $s(w)=\sum_t\epsilon_{t,w(t)}$. Suppose $m$ distinct variables $v_i$ have incident factors $f_i$ with $s(a_{v_i})\le-L$ and $s(d_{f_i})\ge L$. If $c_t=\sum_g c_{t,g}$, then
\[m(L-R)\le\sum_{t=1}^Q c_t.\]
The factors may repeat and the measurement spaces may vary arbitrarily with $t$.
\end{theorem}
\begin{proof}
Let $A(T)_{ij}=P_{f_i,v_i}(T)_{v_j}$. Then $A(0)=I_m$, so $\det A$ is nonzero and has degree at most $mR$. At $t$, let $u_{t,g}$ and $v_{t,g}$ count the selected rows assigned to $g$ and selected columns in $g$. Form $M_t$ by retaining entries of $A(t)$ only where row and column groups differ. In each column group its rows lie in the restricted row space of $C_{t,g}$, so $\rank M_t\le c_t$. The remainder has matching-group blocks, hence
\[\rank A(t)\le c_t+\sum_g\min(u_{t,g},v_{t,g})=m+c_t-\Delta_t,\quad \Delta_t=\tfrac12\sum_g|u_{t,g}-v_{t,g}|.\]
If a polynomial matrix has rank $\rho$ at $t$, constant invertible row operations make $m-\rho$ rows vanish there; factoring $T-t$ from each proves that $(T-t)^{m-\rho}$ divides its determinant. This elementary multiplicity fact is also a consequence of Smith form~\cite{matrixpoly}. Summing over distinct points,
\[mR\ge\sum_t(m-\rank A(t))\ge\sum_t\Delta_t-\sum_t c_t.\]
Finally $2mL\le\sum_{t,g}\epsilon_{t,g}(u_{t,g}-v_{t,g})\le2\sum_t\Delta_t$. Combining the inequalities proves the claim.
\end{proof}

\section{Expanding scopes}
Assume each variable occurs in at most $b$ scopes and every $W\subseteq[n]$ with $n/4\le |W|\le n/2$ is crossed by at least $h|W|$ factors, where a crossing factor meets both $W$ and its complement. Here $b$ and $h>0$ are fixed; these are additional hypotheses, not part of the layout definition. Set $\alpha=h/(8b)$.
\begin{lemma}\label{lem:prob}
For $Q\ge(512/\alpha)^2$, $0\le L\le\alpha Q/64$, and two schedules at Hamming distance at least $\alpha Q$, independent uniform signs $\epsilon_{t,g}$ satisfy
\[\Pr[s(a)\le-L,\ s(d)\ge L]\ge c_\alpha\exp[-C_\alpha(L^2/Q+1)]\]
for positive constants depending only on $\alpha$.
\end{lemma}
\begin{proof}
Let $\ell\ge\alpha Q$ be the distance. At the $\ell$ mismatched times the two signs $\epsilon_{t,a(t)}$ and $\epsilon_{t,d(t)}$ are independent; let $k$ count those where they disagree. Then $k\sim\mathrm{Bin}(\ell,1/2)$ and Paley--Zygmund gives $\Pr(k\ge\ell/4)\ge1/8$. Conditional on these positions, $s(a)=C-Z$ and $s(d)=C+Z$, where $C$ and $Z$ are independent sums of $Q-k$ and $k$ fair signs. Chebyshev gives $\Pr(|C|\le2\sqrt Q)\ge3/4$. Put $u=L+2\sqrt Q$ and $w=u+2\sqrt k$. If $w\le k/4$, exponential tilting with $k\tanh\lambda=w$ and Chebyshev under the tilted law yield
\[\Pr(Z\ge u)\ge\tfrac34\exp\{-3w(u+4\sqrt k)/k\}.\]
Indeed $\lambda\le3w/k$, and the tilted probability of $|Z-w|<2\sqrt k$ is at least $3/4$; undo the tilt using $(\cosh\lambda)^k\ge1$. The interval argument needs no parity adjustment. On $k\ge\alpha Q/4$, the stated range gives $w\le L+4\sqrt Q\le\alpha Q/64+\alpha Q/128<\alpha Q/16\le k/4$. On the event $|C|\le2\sqrt Q$ and $Z\ge u$ we have $s(d)=C+Z\ge L$ and $s(a)=C-Z\le-L$. The exponent is $O_\alpha(L^2/Q+1)$, and multiplying the three probability bounds proves the lemma.
\end{proof}
\begin{theorem}\label{thm:expand}
Under the above expansion hypotheses, every layout satisfies, for a constant $c>0$ depending only on $b,h$,
\[F\ge \frac c{\sqrt Q},\qquad F\ge c\frac{\log_2(2Q)}{\sqrt Q}.\]
If also $Q\log_2 r\le K\log_2 n$ for a fixed $K$, then for sufficiently large $n$,
\[F\ge c_{b,h,K}\frac{(\log_2\log_2 n)^{3/2}}{\sqrt{\log_2 n}}.\]
\end{theorem}
\begin{proof}
For $r\ge8$, partition parts can be greedily united to give $n/4\le|W|\le n/2$. Each crossing factor has an incidence whose factor and variable groups differ, so there are at least $hn/4$ such incidences at each time. There are at most $bn$ incidences, hence at least $hn/8$ have schedule distance at least $\alpha Q$. Lemma~\ref{lem:prob}, expectation, and selection of one incidence per variable give signs and $m\ge cn\exp(-CL^2/Q)$ selected variables. Theorem~\ref{thm:det} and $c_t\le2\tau n$ imply, whenever $L>R$ is in the lemma's range,
\begin{equation}\label{eq:threshold}
 \tau\ge c\frac{L-R}{Q}\exp(-CL^2/Q).
\end{equation}
For $R\le\sqrt Q$, take $L=R+\lceil\sqrt Q\rceil$; for the lemma's lower bound on $Q$, this is admissible and gives $\tau\ge c/\sqrt Q$, hence $F\ge\tau\log_2 3\ge c/\sqrt Q$ since $D\ge1$. For $R>\sqrt Q$, the first term of \eqref{eq:F} is at least $R/Q>1/\sqrt Q$. Bounded $Q$ and $r<8$ are handled by reducing $c$, using respectively $R/Q\ge1/Q$ and $1/r>1/8$.

For the stronger bound, integer measurement rows have support at most $D$. At each time the union of their supports has size at most $2\tau nD$. An off-assigned incidence $(f,v)$ outside this union forces $P_{f,v}(t)_v=0$; its diagonal polynomial is nonzero at zero and has at most $R$ roots. Summing over at most $bn$ incidences gives
\begin{equation}\label{eq:support}
 R/Q+2\tau D\ge h/(4b).
\end{equation}
If $R>\sqrt Q$, concavity of $x\log_2(2/x)$ on $[Q^{-1/2},1]$ gives the claimed bound. Otherwise, for large $Q$, \eqref{eq:support} gives $D\ge h/(16b\tau)$ and the preceding argument gives $\tau\ge c_0/\sqrt Q$. The function $\tau\log_2(1+z/\tau)$ is increasing for $z>0$, so the measurement term is $\Omega(\log(2Q)/\sqrt Q)$. Again bounded $Q$ and $r<8$ are absorbed into the constant.

For the capacity conclusion put $N=\log_2 n$. If $r\le N^{1/4}$, then $1/r\ge N^{-1/4}$ dominates the claimed bound. Otherwise $Q\le4KN/\log_2 N$. The function $\log_2(2Q)/\sqrt Q$ decreases for sufficiently large $Q$; evaluation at that upper endpoint proves the result, while bounded $Q$ gives a constant.
\end{proof}

\section{Auxiliary coordinates}
Let $B$ be a nonempty finite label set. For each $u\in B$ let $a_u,d_u\in[r]^Q$, and let $P_u(T)\in\mathbb Q[T]^B$ have degree at most $R$ and $P_u(0)=e_u$. At time $t$ let $\Lambda_t\subseteq\mathbb Q^B$ have dimension at most $\varepsilon|B|$, and suppose $P_u(t)-\lambda_{u,t}$ is supported on $\{w:a_w(t)=d_u(t)\}$ for some $\lambda_{u,t}\in\Lambda_t$. Corrections may depend on both $u$ and $t$.
\begin{theorem}\label{thm:aux}
Fix $\alpha,\beta,K>0$. If at least $\beta|B|$ labels satisfy $\dist(a_u,d_u)\ge\alpha Q$ and $0<\varepsilon\le K/Q$, then, for sufficiently large $Q$,
\[R\ge c\min\{Q,\sqrt{Q\ln(1/\varepsilon)}\},\]
where $c>0$ depends only on $\alpha,\beta,K$. In particular this holds on the full label space $B=([r]\times[r])^Q$ for $r\ge2$, with constants independent of $r$.
\end{theorem}
\begin{proof}
Lemma~\ref{lem:prob} selects $m\ge c|B|e^{-CL^2/Q}$ long labels. The determinant proof applies with the whole correction matrix of rank at most $\varepsilon|B|$ at each time, yielding
\[\varepsilon\ge c(L-R)Q^{-1}e^{-CL^2/Q}.\]
If $R<\sqrt Q$, taking $L=R+\lceil\sqrt Q\rceil$ contradicts $\varepsilon\le K/Q$ for large $Q$. If $2R>\alpha Q/64$, the linear branch of the conclusion holds. Otherwise take $L=2R$: the inequality gives $\varepsilon\ge(c/\sqrt Q)e^{-4CR^2/Q}$. As $\ln(1/\varepsilon)\ge\ln Q-\ln K$, rearrangement gives $R^2\ge cQ\ln(1/\varepsilon)$ for large $Q$. On the full space the mean distance is $(1-1/r)Q\ge Q/2$, so at least one third of labels have distance at least $Q/4$.
\end{proof}
The long-pair assumption cannot be dropped: even with balanced schedules, take $B=[r]^Q$, $a_u=d_u=u$, $P_u=e_u$, $R=1$, and one-dimensional correction spaces (or zero corrections). Then $\varepsilon=1/|B|$ is allowed and no such degree bound follows. In OpenAI's incidence embedding, $\tau=rb\varepsilon$ is a sufficient physical budget, not necessarily the smallest possible budget; the auxiliary theorem itself incurs no $\ln(rb)$ loss.

\section{Partner times}
Let $A(T)$ be an $m\times m$ polynomial matrix of degree at most $R$ with $A(0)=I$. All evaluation times are distinct and nonzero. At time $t$, rows and columns have group labels $d_i(t),a_j(t)$. Let $N_t=\{(i,j):d_i(t)=a_j(t)\}$ be the unmeasured entries, and let $I_t=\{i:d_i(t)\ne a_i(t)\}$. Let $\mathrm{cost}_t$ be the sum, over column groups, of the ranks of the measured (off-$N_t$) column blocks. This is at most the number of layout measurement rows when $A$ is a selected incidence matrix.
\begin{proposition}\label{prop:partner}
For each $t$ in a set $\mathcal T$, suppose there is a set $P(t)$ of $R$ distinct nonzero times, excluding $t$, at which every entry of $N_t$ is measured. For $t'\in P(t)$ let $\beta(t,t')$ be the maximum, over column groups $g'$ at $t'$, of the number of labels $h$ for which some column has $a_j(t)=h$ and $a_j(t')=g'$. Put
\[\lambda=\max_{t'}\sum_{t\in\mathcal T:t'\in P(t)}\beta(t,t').\]
Then, with cost summed over all evaluation times,
\[(1+\lambda)\sum_t\mathrm{cost}_t\ge\sum_{t\in\mathcal T}|I_t|-mR.\]
\end{proposition}
\begin{proof}
Let $G_t(T)$ retain the entries of $A(T)$ in $N_t$ with row index in $I_t$, and zero all other entries. These entries are off-diagonal, so $G_t(0)=0$ and $\deg G_t\le R$. At $t$, the measured part has rank at most $\mathrm{cost}_t$, while the unmeasured part in rows outside $I_t$ has rank at most $m-|I_t|$. Thus the corank of $A(t)$ is at least $|I_t|-\mathrm{cost}_t-\rank G_t(t)$.

At a partner $t'$, within each column group $g'$, $G_t(t')$ splits into at most $\beta(t,t')$ rectangles, indexed by $h$, with rows $i\in I_t$, $d_i(t)=h$, and columns $a_j(t)=h$, $a_j(t')=g'$. Each is a row-and-column restriction of the measured block at $t'$. Hence $\rank G_t(t')\le\beta(t,t')\mathrm{cost}_{t'}$. Lagrange interpolation through $0$ and $P(t)$ gives $\rank G_t(t)\le\sum_{t'\in P(t)}\beta(t,t')\mathrm{cost}_{t'}$. Summing the corank inequalities and using the determinant multiplicity bound $\sum_t\mathrm{corank}\,A(t)\le mR$ gives the result.
\end{proof}
\begin{corollary}\label{cor:coset}
Let indices be $x\in\mathbb F_2^k$, $m=2^k$. For each of $\Delta$ nonzero directions $y_\ell$, take one antisymmetric time with $a_x=\langle x,y_\ell\rangle$, $d_x=1-a_x$, and $R$ symmetric times with $d_x=a_x$ using the same direction. Then every such $A$ has total cost at least $m(\Delta-R)$; for $R\le\Delta$ this is attained.
\end{corollary}
\begin{proof}
At an antisymmetric time the measured part is block diagonal and the unmeasured part is block antidiagonal. Write $G_y$ for the block-antidiagonal projection. The determinant bound gives $\sum_{t\,\mathrm{anti}}\mathrm{cost}_t\ge m(\Delta-R)-\sum_{t\,\mathrm{anti}}\rank G_{y_t}(A(t))$. Each $G_y(A(T))$ vanishes at zero; interpolation through its $R$ symmetric partners bounds its rank at the antisymmetric time by the sum of their measured-block ranks. The partner sets are disjoint, giving the bound. For $R\le\Delta$, choose a scalar polynomial $p$ of degree $R$, with $p(0)=1$ and roots at any $R$ antisymmetric times, and take $A(T)=p(T)I$: its total cost is exactly $m(\Delta-R)$. At symmetric times the unmeasured part is block diagonal and contains the identity. Consequently the determinant rank estimate of Theorem~\ref{thm:det} alone charges nothing there; interpolation brings those times into the argument.
\end{proof}
In the family of Corollary~\ref{cor:coset}, a row $x$ and a column $x'$ differ at the antisymmetric time of direction $y_\ell$ when $\langle x,y_\ell\rangle=\langle x',y_\ell\rangle$ and at its $R$ symmetric times otherwise, so $\dist(d_x,a_{x'})\ge\Delta$ for every pair. The first public version of this paper (Zenodo, version 1.0) conjectured that separation alone suffices: there is an absolute $c>0$ such that, for arbitrary row and column schedules with $\dist(d_i,a_j)\ge\Delta>R$ for all $i\le j$, every degree-$R$ matrix $A$ with $A(0)=I$ has total measured-block rank at least $c\,m(\Delta-R)$. In Tulsyan's second theorem, ARV~\cite{arv} separates two large sets of variable schedules, and linkedness supplies many incidence labels with all selected row--column pairs far apart, so the all-pairs case of this bound would have extended that route to time-dependent measurements. The next section shows that the conjecture is false, already in the all-pairs case with two groups.

\section{Partner times cannot be dropped}
The \emph{paired $j$-cube} has indices $x\in\F_2^j$, $m=2^j$, two groups, and times $[Q]$ with $Q=2j$. Direction $\ell\in\{0,\dots,j-1\}$ has a symmetric time $\sigma_\ell=2\ell+1$, with $a_x=d_x=x_\ell$, and an antisymmetric time $\alpha_\ell=2\ell+2$, with $a_x=x_\ell$, $d_x=1-x_\ell$. This is the family of Corollary~\ref{cor:coset} with $y_\ell=e_\ell$ and $\Delta=j$, but with one symmetric time per direction instead of $R$. The entry $(x,y)$ is measured at $\alpha_\ell$ if and only if $x_\ell=y_\ell$, and at $\sigma_\ell$ if and only if $x_\ell\ne y_\ell$. Hence $\dist(d_x,a_y)=j$ for every pair $x,y$, including $x=y$.
\begin{theorem}\label{thm:far}
Let $q\ge3$ and $s\ge1$, and put $j=s(2q^2+2q+2)$, $R=2sq(q+1)$ and $\Delta=j$, so that $\Delta-R=2s$. On the paired $j$-cube there is a matrix $A(T)$ of degree at most $R$, with $A(0)=I$ and entries in a number field, whose total measured-block rank is
\[(q+1)\,4^{-q}\,m(\Delta-R).\]
Repeating every row and column schedule $e$ times gives a rational matrix with the same ratio. Hence the far-pair bound fails for every $c>0$, over $\mathbb Q$, with two balanced groups, times $[Q]$, and $\dist(d_i,a_j)=\Delta$ for all pairs.
\end{theorem}
The matrices are built from Clifford algebras. Fix $q$ directions $P$. For $w\in\F_2^q$ put $h_w(T)=\prod_{p\in P}(T-\rho_p(w))$, where $\rho_p(w)=\sigma_p$ if $w_p=1$ and $\rho_p(w)=\alpha_p$ otherwise. For $(w,v)\in\F_2^q\times\F_2^q$ let $G_{w,v}$ be the $2^q\times2^q$ signed permutation matrix with entry $(-1)^{v\cdot y}$ at $(y+w,y)$. Then $G_{w,v}^2=(-1)^{w\cdot v}I$, $G_{w,v}$ has trace zero unless $(w,v)=0$, and $G_{w,v}$, $G_{w',v'}$ anticommute if and only if $w\cdot v'+w'\cdot v=1$. A \emph{$q$-family} is a list of $2q+1$ pairwise anticommuting labels $(w,v)$ with distinct $w$, such that the $2q+1$ polynomials $h_w^2$ are linearly independent, hence a basis of the polynomials of degree at most $2q$.
\begin{lemma}\label{lem:cluster}
Given a $q$-family on $P$ and distinct nonzero times $t_1,\dots,t_{2q+2}$ outside $\{\sigma_p,\alpha_p:p\in P\}$, there is a $2^q\times2^q$ matrix $M(T)$ of degree at most $q+1$ over a number field with $M(0)=I$, such that every entry measured at a time $\sigma_p$ or $\alpha_p$, $p\in P$, vanishes there, $\det M=g^{2^{q-1}}$ with $g=\prod_l(1-T/t_l)$, and $\rank M(t_l)=2^{q-1}$ for every $l$.
\end{lemma}
\begin{proof}
Let $F=\prod_{p\in P}(1-T/\alpha_p)$ and $H_w=Th_w$. Choose $b$ so that $F^2(1+bT)^2-g$ has no linear term. That difference is then $-T^2$ times a polynomial of degree at most $2q$, so the basis property gives unique $k_w$ with $g=F^2(1+bT)^2-\sum_wk_wH_w^2$. Choose $c_w$ with $c_w^2=(-1)^{w\cdot v}k_w$, and put $d=F(1+bT)$, $N=\sum_wc_wH_wG_{w,v}$ and $M=dI+N$. Then $M(0)=I$ and $\deg M\le q+1$. Because the $w$ are distinct, the entry $(x,y)$ of $N$ involves only the term with $w=x+y$. If it is measured at $\alpha_p$, then $w_p=0$, and both $H_w$ and $d$ (through $F$) vanish at $\alpha_p$. If it is measured at $\sigma_p$, then $w_p=1$ and $x\ne y$, and $H_w$ vanishes at $\sigma_p$.

Anticommutation gives $N^2=\kappa I$ with $\kappa=\sum_wk_wH_w^2=d^2-g$. The polynomial $\kappa$ is nonzero, since otherwise $g=d^2$ would vanish at $\alpha_p$. Over an algebraic closure of the field of rational functions, $N$ is diagonalizable with eigenvalues $\pm\sqrt\kappa$, and both multiplicities are $2^{q-1}$ because $N$ has trace zero. Hence $\det M=g^{2^{q-1}}$. At a root $t_l$ of $g$ we have $\kappa(t_l)=d(t_l)^2$. If $d(t_l)\ne0$, then $M(t_l)^2=2d(t_l)M(t_l)$, so $M(t_l)/(2d(t_l))$ is an idempotent of trace $2^{q-1}$. If $d(t_l)=0$, then $M(t_l)^2=0$, so $\rank M(t_l)\le2^{q-1}$, while the nullity is at most the order $2^{q-1}$ of the root of $\det M$, since $t_l$ is a simple root of $g$.
\end{proof}
\begin{lemma}\label{lem:families}
For every $q\ge3$ there is a $q$-family on any $q$ consecutive directions.
\end{lemma}
\begin{proof}
Translating all times by a constant replaces each $h_w(T)$ by $h_w(T-c)$, a unipotent change of basis, so linear independence does not depend on the placement. For $q=3,4,5$ the families
\begin{align*}
&(1,0),(2,1),(3,5),(4,3),(5,1),(6,3),(7,5),\\
&(0,3),(1,0),(2,1),(5,1),(6,11),(9,5),(10,3),(13,5),(14,11),\\
&(5,26),(6,18),(7,23),(10,30),(14,17),(16,0),(17,31),(19,19),(22,27),(24,22),(28,17)
\end{align*}
(bit $i$ of each integer is direction $i$) satisfy the axioms, and the coefficient matrices of their polynomials $h_w^2$ have determinants $4096$, $715456512$ and $176287042228935720960$ in exact arithmetic. For larger $q$, let $\Gamma_1$ be a $q_1$-family on $P_1$ and $\Gamma_2$ a $q_2$-family on disjoint $P_2$, and fix $\gamma^*\in\Gamma_1$. The labels $\gamma\otimes I$ for $\gamma\in\Gamma_1\setminus\{\gamma^*\}$ and $\gamma^*\otimes\delta$ for $\delta\in\Gamma_2$ number $2(q_1+q_2)+1$. They anticommute pairwise, because the symplectic form adds over tensor factors, and their $X$-parts $(w_\gamma,0)$ and $(w^*,w_\delta)$ are distinct. Their polynomials are $h^{(1)2}_{w_\gamma}F_2'^2$ and $e\,h^{(2)2}_{w_\delta}$, where $F_2'=\prod_{p\in P_2}(T-\alpha_p)$ and $e=h^{(1)2}_{w^*}$. A linear relation reads $AF_2'^2=-eB$, with $A$ in the span of the first kind and $B$ in the span of the second. The roots of $e$ are times of $P_1$ and those of $F_2'$ are times of $P_2$, so $e$ divides $A$. Since $\deg A\le2q_1=\deg e$, $A=\lambda e$, and independence of $\Gamma_1$ forces $\lambda=0$ and all coefficients of $A$ to vanish. Then $B=0$, and independence of $\Gamma_2$ finishes the proof. Every $q\ge6$ is a sum of $3$s and $4$s.
\end{proof}
\begin{proof}[Proof of Theorem~\ref{thm:far}]
Put $h=2q+2$ and $k=2q$. Split the directions into $s$ groups, each consisting of $k$ blocks of $q$ consecutive directions and $h$ \emph{root directions}. Within a group, let every block carry the matrix of Lemma~\ref{lem:cluster} (with Lemma~\ref{lem:families}) whose roots are the times $\alpha_\ell$ of that group's root directions. Let $A$ be the Kronecker product of these $sk$ matrices and of the identity on the root bits. Then $A(0)=I$ and $\deg A\le sk(q+1)=R$. An entry of $A$ measured at a time of a block direction is a product that contains a measured entry of that block's factor, so it is zero. At $\sigma_\ell$ for a root direction $\ell$, the measured entries have $x_\ell\ne y_\ell$, where $A$ vanishes. At $\alpha_\ell$ the measured part consists of the two diagonal blocks $x_\ell=y_\ell$, so the cost there is $\rank A(\alpha_\ell)=m\prod_i\rank M_i(\alpha_\ell)/2^q$. The $k$ factors of the group containing $\ell$ have rank $2^{q-1}$ there, and all other factors are invertible. The total cost is therefore $m\cdot sh\cdot2^{-k}$, while $\Delta-R=s(kq+h)-sk(q+1)=2s$, which gives the ratio $(q+1)4^{-q}$.

For the rational version, let $K$ be the number field generated by the entries of $A$, of degree $e$ over $\mathbb Q$, and replace each entry by its $e\times e$ regular representation. Copy every row and column schedule $e$ times. Measured index sets are unions of whole copy-blocks, the regular representation of a $K$-matrix has $\mathbb Q$-rank $e$ times its $K$-rank, and $m$ also grows by the factor $e$. Distances and degree are unchanged.
\end{proof}
The construction is optimal within its own class. With $c_\ell$ factors sharing the root direction $\ell$, convexity of $2^{-c}$ shows that $\sum_\ell2^{-c_\ell}\ge(q+1)4^{-q}(\Delta-R)$ for factors of size at most $q$. In terms of the degree alone ($s=1$), its cost is $m(\Delta-R)\,2^{-(1+o(1))\sqrt{2R}}$. The smallest instance that beats $m(\Delta-R)$ is much smaller: on the paired $4$-cube with $R=2$, an explicit real matrix has cost $24<32$. This is consistent with our earlier searches, which stopped at $m=12$. The factorization $(dI+N)(dI-N)=gI$ is the norm form of a Clifford algebra, as in matrix factorizations~\cite{eisenbud}. Theorem~\ref{thm:far} concerns the abstract matrix statement only. It does not construct a layout: it says nothing about coefficient sizes, measurement-row norms, or linked expanding scopes, and it leaves open whether those, or enough partner times as in Corollary~\ref{cor:coset}, restore a bound of order $m(\Delta-R)$.

\section{Limitations and verification}
It remains open whether, on the linked bounded-incidence scope systems used in Tulsyan's second theorem, sufficiently small time-dependent measurement budgets force $R=\Omega(Q/\sqrt{\log n})$ without a capacity assumption. A simple equal-marginal example has row schedules $00,11$ and column schedules $01,10$: every coordinate imbalance is zero although all cross-pairs are separated. It shows why the additive determinant charge does not directly capture the nonlinear separation used in the fixed-correction argument.

Even with the capacity condition, our bound is smaller by a factor $\sqrt{\log\log n}$ than the order obtained in \cite{openai}, namely $(\log\log n)^2/\sqrt{\log n}$. This cannot be closed by optimizing \eqref{eq:threshold} and combining it only with \eqref{eq:support}. Put $N=\log_2n$ and take
\[Q=\lfloor N/\log_2N\rfloor,\quad R=1,\quad r=\lceil\sqrt Q\rceil,\quad \tau=\kappa_1/\sqrt Q,\quad D=\lceil \kappa_2\sqrt Q\rceil,\]
with sufficiently large constants $\kappa_1,\kappa_2$. These formal parameters satisfy the displayed necessary inequalities and $Q\log_2r=(1/2+o(1))N$, while $F=\Theta(\log Q/\sqrt Q)$. They do not construct a layout: additional polynomial and incidence constraints may exclude them. The capacity condition itself is not part of the layout definition.

The repository's exact certificate checker verifies finite layouts and determinant certificates, including equality on a four-cycle and genuinely varying measurement spaces on the four-vertex complete graph. It checks a fixed-schedule feasibility claim and small sign-mixture probabilities; the latter do not establish the asymptotic probability constants. The partner checker verifies a projection identity and the corresponding block-rank identity exactly. The far-pair checker verifies the base families of Lemma~\ref{lem:families}, the recursive families up to $q=10$, and the matrices of Lemma~\ref{lem:cluster} for $q=3,4,5$, all in exact arithmetic. These finite calculations do not replace the general proofs.

\section*{Acknowledgment}
This work was carried out with AI research assistants. The proofs were cross-checked by an independent AI reviewer and the finite claims by exact computation.
\begin{thebibliography}{9}
\bibitem{openai} OpenAI. \emph{The Exponential Time Hypothesis is False}. October 17, 2026; revised October 25, 2026. \href{https://github.com/openai/math/tree/main/preprints/The-Exponential-Time-Hypothesis-is-False-October-25-2026}{Preprint and code}.
\bibitem{anthropic} Claude (Anthropic). \emph{A $2^{o(n)}$ Time Algorithm for 3-SAT}. October 26, 2026. \href{https://github.com/anthropics/formal-math/tree/main/eth}{Preprint and code}.
\bibitem{tulsyan} Vaibhav Tulsyan. \emph{The parity barrier in OpenAI's Engine I}. October 29, 2026. \href{https://github.com/xennygrimmato/eth-hamiltonicity/blob/a4f924c1419cbc43564c3792925d8bb182ca932c/research/permanent/parity-barrier.md}{Research note}.
\bibitem{olivier} Antoine Olivier. \emph{openai-eth-section4-lower-bound.md}. October 28, 2026. \href{https://gist.github.com/antoine-olivier/d269b3b449d3a2cdb021dc0ee5d0fc82}{Research note}.
\bibitem{matrixpoly} Israel Gohberg, Peter Lancaster, and Leiba Rodman. \emph{Matrix Polynomials}. Academic Press, 1982.
\bibitem{eisenbud} David Eisenbud. Homological algebra on a complete intersection, with an application to group representations. \emph{Transactions of the AMS} 260(1), 1980.
\bibitem{arv} Sanjeev Arora, Satish Rao, and Umesh Vazirani. Expander flows, geometric embeddings and graph partitioning. \emph{Journal of the ACM} 56(2), 2009.
\end{thebibliography}
\end{document}
